\documentclass{beamer}
%[aspectratio=169]   \usepackage[czech]{babel}
\usepackage{apo-lecture-en}
\usepackage{pdfpages}
\usepackage{pdfcomment}
\usepackage{listings}
\usepackage{array,multirow}
\usepackage{environ}

\subtitle{Lecture 13. Virtual Machines and Virtualization}
\author{Pavel Píša \phantom{xxxxxxxxx} Petr Štěpán \\ \small\texttt{pisa@fel.cvut.cz}\phantom{xxxx}\small\texttt{stepan@fel.cvut.cz}}

\begin{document}

\maketitle

\section{Multilevel Computer Organization}

\begin{frame}
\frametitle{Today's Lecture Objective}

\begin{itemize}
 \item Multilevel Computer Organization
 \item Virtualization
\end{itemize}
\end{frame}

\begin{frame}
\frametitle{Multilevel Computer Organization}

\begin{itemize}
 \item Machine code (language)
 \begin{itemize}
  \item instruction set -- specific for given architecture
  \item language L1 level -- alphabet {0,1}
  \item it is necessary to convert the program into it for native execution
  \item for application design by humans and generators, higher-level languages ​​L2 + others are more suitable
 \end{itemize}
 \item Executing a program in L2 on a machine that provides L1
 \begin{itemize}
  \item Compilation -- instructions in L2 are replaced by a sequence of instructions in L1
  \item Interpretation -- a program in L1 processes a program in L2 as data (slower) jako data
 \end{itemize}
 \item Development of machine levels
 \begin{itemize}
  \item first computers -- ordinary machine level - 1st level
  \item 50s -- Wilkes -- microprogram. - 2nd level
  \item 60s -- operating systems - 3rd level
  \item compilers, programming languages ​​- 4th level
  \item user application programs - 5th level
  \item problem-oriented languages ​​and environments
 \end{itemize}
\end{itemize}

\end{frame}

\begin{frame}
\frametitle{The Virtual Computer}

Hardware (HW) and software (SW) are logically equivalent (can be replaced with each other) the battle and merging of RISC and CISC processors
\vspace{0.2cm}
At level $i$ we introduce a virtual computer $M_i$ with language $L_i$.
The program in $L_i$ is translated or interpreted by computer $M_{i-1}$ etc.

\begin{center}
\includegraphics[width=0.75\textwidth]{generated/computer/mutilevel-machine-en.pdf}
\end{center}

\end{frame}

\begin{frame}
\frametitle{Current Multi-level Computer Architecture}

\begin{center}
\includegraphics[width=0.75\textwidth]{generated/computer/mutilevel-mcode-to-highlev-en.pdf}
\end{center}

\end{frame}

\begin{frame}
\frametitle{Processes and Their States}

\begin{description}
 \item[process] a running program (program -- passive, process -- active)
 \item[process state] information that allows a process to be restarted when it is stopped
 \begin{enumerate}
  \item program code (.text)
  \item next. instructions (PC/IP)
  \item variable values ​​and data
  \item states and positions of all I/O (positions in open files, etc.)
  \item address space layout/occupation (memory context)
 \end{enumerate}
 \item[assumption] process P does not change its program itself, if so, then the changed parts are part of the state, data
 \item[state vector] variables of the process state component -- they are changed by the program, OS or HW
   $$PROCESS = PROGRAM + STATE \medskip VECTOR$$
\end{description}
\end{frame}

\begin{frame}
\frametitle{Conventional Machine Level (ISA) (M2)}

Defined by the Instruction Set (often referred to as processor architecture) ISA -- Instruction Set Architecture $\times$ Microarchitecture (Implementation in M1 or directly in the logical design of the processor)

\begin{itemize}
 \item structure of the computer, processor, I/O channels, buses, organization and accesses to memory, register organization, etc.
 \item instruction set, data format and storage in memory, addressing, notebook memory (registries)
\end{itemize}

Instruction format

\begin{itemize}
 \item consists of fields stored in one or more consecutive words (sometimes bytes granularity)
 \item operation character (opcode)
 \item immediate operands, constants embedded in the instruction
 \item address parts can be memory addresses or references, register numbers and these can be used either directly or as memory addresses (indirect addressing)
\end{itemize}
\end{frame}

\section{Virtualization}

\begin{frame}
\frametitle{Virtualization}

\begin{itemize}
 \item Virtualization hides the implementation/properties of lower layers (reality) and presents an environment with the required properties
 \item In computer technology, we divide virtualization into
 \begin{itemize}
  \item Purely application/at the language and compiled code (byte-code) level – virtual machines, e.g. JVM, dotNET
  \item Emulation and simulation typically different computer architectures (also called cross-virtualization)
  \item Native virtualization - isolated environment providing identical type of architecture for unmodified OS
  \item Virtualization with full support directly in HW
  \item Partial virtualization -- typically only address spaces
  \item Paravirtualization -- system must be modified to run in the offered environment
  \item OS-level virtualization -- only separate user environments (containers)
 \end{itemize}
\end{itemize}
\end{frame}

\begin{frame}
\frametitle{Operating System Level and Isolation of Applications}

\begin{itemize}
 \item Threads can also be considered as virtualization of the processor state and processes as virtualization of the entire environment including memory separation
 \item Virtual machines and interpreters within applications (browser, Python in LibreOffice, etc.)
 \begin{itemize}
  \item direct interpretation in text form is rare, mostly translation into binary representation (bytecode) and from there often into native code (JVM)
 \end{itemize}
 \item Containers
 \begin{itemize}
  \item chroot -- separate view of the directory structure (Mount -- mnt)
  \item LXC (Linux Containers) -- separability of various kernel namespaces (namespace) and resource prioritization (cgroups)
  \begin{itemize}
   \item Mount (mnt), Process ID (pid), Network (net), Inter-process Communication (ipc), UTS (host+domain), User ID (user), Control group (cgroup) Namespace, Time Namespace
   \item see system call manual: \texttt{unshare(2)}, \texttt{clone(2)}
  \end{itemize}
  \item Docker -- mostly on top of LXC
  \item systemd-nspawn
  \item Solaris Containers (Zones)
  \item FreeBSD jail
 \end{itemize}
\end{itemize}
\end{frame}

\begin{frame}
\frametitle{Virtualization at the Whole Machine Level}

\begin{itemize}
 \item Host computer (Host, Domain DOM 0)
 \item Guest system (Guest)
 \item Processor for the guest system
 \begin{itemize}
  \item for the native case, regular code/unprivileged instructions processed directly by the CPU
  \item for the cross-architectural, interpretation of instructions by the emulator (program in DOM 0), or acceleration
 \end{itemize}
 \item Privileged instructions in the guest system
 \begin{itemize}
  \item causes exceptions that are handled by the monitor/hypervisor, emulation
  \item CPU has support for HW virtualization (AMD-V, Intel VT-x), shadow paging tables
 \end{itemize}
 \item Peripherals/devices
 \begin{itemize}
  \item IO access and memory-mapped peripherals cause exceptions and the hypervisor simulates their operation
  \item the guest system adapts itself so that passed the requests directly in a format that the hypervisor understands (custom drivers, etc.)
 \end{itemize}
\end{itemize}
\end{frame}

\begin{frame}
\frametitle{Hypervisor}

\begin{itemize}
 \item ensures the starting and stopping of domains
 \item monitors their activity and handles exceptions - mainly emulates the activity of privileged instructions
 \item ensures the distribution of memory and computing power to guest systems
 \item emulates the activity of peripherals and passes data to physical device and network drivers at the host system level
 \item can be implemented
 \begin{itemize}
  \item in the user space of the host system (QEMU)
  \item using support in HW and the OS kernel (KVM)
  \item as a separate system/microkernel that uses the system in one domain (DOM 0) to communicate with physical devices and forwards data from this system to others (XEN) 
 \end{itemize}
\end{itemize}
\end{frame}

\begin{frame}
\frametitle{Debian Packages for Virtualization}

\begin{itemize}
 \item Xen
 \item KVM
 \begin{itemize}
  \item Allows full GPU propagation via PCIe passthrough
  \item VirtIO
  \item VirGL
  \item KVMGT (Intel Haswell-based GPUs for now)
  \item virtfs
  \item virtio-net
 \end{itemize}
 \item Qemu System -- mainly local and cross-virtualization
 \item qemu-user-static/qemu-user -- userspace only emulation (binfmt\_misc)
\end{itemize}
\end{frame}

\begin{frame}
\frametitle{User Level Extensions and Solutions}

\begin{itemize}
 \item VirtualBox
 \item ganeti -- cluster manager, KVM or Xen.
 \item XCP (or XenServer) manage virtual machines (templates virtual storage and networking) 
 \item OpenStack (which uses KVM or Xen) large scale 
 \item libvirt -- KVM, Xen and more (CLI virsh, API) 
 \item virt-manager - above libvirt 
 \item kpartx - convenient utility for accessing partitions in an image file or LV from a virtual machine 
 \item xmount - tool to crossmount between multiple input and output hard disk images 
 \item SystemBuildTools includes some tools to build VMs
\end{itemize}
\end{frame}

\begin{frame}
\frametitle{Packages for Container-based Solutions}

\begin{itemize} 
 \item nspawn - containers provided by systemd. 
 \item Incus, LXD and LXC
 \item Docker and Podman
 \item OpenVz (DEPRECATED)
\end{itemize}

Partial isolation only

\begin{itemize}
 \item UserModeLinux
 \item chroot
 \item Schroot (based on chroot)
 \item Bubblewrap
 \item Isolate
\end{itemize}
\end{frame}

\end{document}

